
\section{VI - Reconnaissance de mot}
\begin{frame}{VI - Reconnaissance de mot}
Un moyen simple de reconnaitre un mot est d'en faire le spectrogramme puis de résoudre le problème de reconnaissance comme si cela était une image. \\
Le site du gouvernement recense ces mots clés pour les appels d'urgences : \\
\begin{enumerate}
  \item[15] malaise, hémorragie, brûlure, intoxication
  \item[17] violences, agression, vol, cambriolage
  \item[18] incendie, gaz, effondrement, électrocution
\end{enumerate}
\begin{figure}
	\begin{subfigure}[]{0.3\textwidth}
		\includegraphics[width=90px]{1-Incendie.jpg}
  		\caption{INCENDIE}
	\end{subfigure}
	\begin{subfigure}[]{0.3\textwidth}
		\includegraphics[width=90px]{1-Intoxication.jpg}
  		\caption{INTOXICATION}
	\end{subfigure}
	\begin{subfigure}[]{0.3\textwidth}
		\includegraphics[width=90px]{1-Cambriolage.jpg}
		\caption{CAMBRIOLAGE}
	\end{subfigure}
\end{figure}
\end{frame}

\section{VII - Reconnaissance vocale}
\begin{frame}{VII - Reconnaissance vocale}
Découpage en élément lexicaux, puis décodage du mot.
\begin{figure}
	\begin{subfigure}[]{0.3\textwidth}
		\includegraphics[width=90px]{1-A.jpg}
  		\caption{A}
	\end{subfigure}
	\begin{subfigure}[]{0.3\textwidth}
		\includegraphics[width=90px]{1-E.jpg}
  		\caption{E}
	\end{subfigure}
	\begin{subfigure}[]{0.3\textwidth}
		\includegraphics[width=90px]{1-I.jpg}
		\caption{I}
	\end{subfigure}
\end{figure}
\begin{figure}
	\begin{subfigure}[]{0.3\textwidth}
		\includegraphics[width=90px]{1-O.jpg}
  		\caption{O}
	\end{subfigure}
	\begin{subfigure}[]{0.3\textwidth}
		\includegraphics[width=90px]{1-U.jpg}
  		\caption{U}
	\end{subfigure}
	\begin{subfigure}[]{0.3\textwidth}
		\includegraphics[width=90px]{2-Matching.png}
	\end{subfigure}
\end{figure}
\end{frame}